%% ============================================================================
%% Section: Evaluation
%% Owns Table tab:datasets, Table tab:coverage, Fig fig:latency, Fig fig:incr.
%% Design is final; every number is a \TODO until bench/run_all.sh (CODE-PLAN
%% WP5) produces paper/numbers.tex. Do not type numbers by hand.
%% ============================================================================
\section{Experimental Evaluation}\label{sec:eval}

In this section, we present an experimental evaluation to assess the
efficiency, the cost, and the correctness of Causal4j. We ask five questions.
\textbf{RQ1 (query performance):} how do structural and identification
queries compare in latency and throughput against in-DBMS and external
alternatives? \textbf{RQ2 (index cost):} what does the index cost to build
and to store as a function of graph size and density, and where is the
crossover with the layer-pruned branch? \textbf{RQ3 (maintenance):} when does
incremental maintenance beat full recompilation, and what does a mixed
workload of edits and queries observe? \textbf{RQ4 (scalability):} how does
latency grow with $|V|$? \textbf{RQ5 (correctness):} do the answers agree
with reference causal tools? All the material necessary to reproduce the
results, together with the hardware specifications, is available on our
online repository\footnote{\TODO{URL of the artifact.}}.

\noindent\textbf{Settings.} \TODO{Hardware (CPU, RAM), JVM version and heap,
Neo4j version (2025.x, embedded through the test harness and as a server for
the Cypher and GDS baselines), JMH configuration: 5 forks, 5 warm-up and 10
measurement iterations per operator; medians with interquartile ranges are
reported.} The index and the baseline share every line of code outside the
representation of the snapshot, so a difference in latency is attributable
to the representation and the algorithms of Section~\ref{sec:index} alone.

\noindent\textbf{Datasets.} We use three families of DAGs
(Table~\ref{tab:datasets}): (i) the reference Bayesian networks of the bnlearn
repository~\cite{scutari2010bnlearn}, taken as causal DAGs, from
\textsc{asia} to \textsc{munin}, which are recognized and reproducible and
span three orders of magnitude in $|V|$; (ii) synthetic DAGs generated by
orienting Erd\H{o}s-R\'enyi and preferential-attachment graphs along a
random vertex order, with $|V|$ from $10^2$ to $10^4$ and average degree in
$\{2, 4, 8\}$ controlled independently, five seeds each, which let us sweep
size and density separately; and (iii) the COVID-19 DAG of the running
example~\cite{steiger2021covid} with $25$ variables and $78$ edges. To
control the workload, every system receives the same $200$ random
treatment-outcome pairs $(T,Y)$ with $Y \in \de(T)$ and the same $200$
random triples $(X,Y,Z)$ with $|Z| \in \{0, 1, 3, 10\}$ per dataset.

\begin{table}[t]\centering
\caption{Datasets. \NOTE{Verify $|V|$ and $|E|$ against the bnlearn
repository before submission; closure memory is $2|V|\lceil|V|/64\rceil$
words.}}
\label{tab:datasets}
\footnotesize
\setlength{\tabcolsep}{3pt}
\begin{tabular}{@{}lrrrr@{}}
\toprule
DAG & $|V|$ & $|E|$ & longest path & closures \\
\midrule
\textsc{asia}       &    8 &    8 & \TODO{} & 128~B \\
\textsc{sachs}      &   11 &   17 & \TODO{} & 176~B \\
\textsc{child}      &   20 &   25 & \TODO{} & 320~B \\
\textsc{alarm}      &   37 &   46 & \TODO{} & 592~B \\
\textsc{hepar2}     &   70 &  123 & \TODO{} & 2.2~KB \\
\textsc{win95pts}   &   76 &  112 & \TODO{} & 2.4~KB \\
\textsc{pathfinder} &  109 &  195 & \TODO{} & 3.5~KB \\
\textsc{andes}      &  223 &  338 & \TODO{} & 14~KB \\
\textsc{pigs}       &  441 &  592 & \TODO{} & 49~KB \\
\textsc{link}       &  724 & 1125 & \TODO{} & 139~KB \\
\textsc{munin}      & 1041 & 1397 & \TODO{} & 283~KB \\
COVID-19            &   25 &   78 & \TODO{} & 400~B \\
synthetic           & $10^2$--$10^4$ & $\bar d{\in}\{2,4,8\}$ & & $\le$25~MB \\
\bottomrule
\end{tabular}
\end{table}

\noindent\textbf{Baselines.} We compare five points along the ``how close to
the DBMS'' axis: (a) the un-indexed baseline of Table~\ref{tab:complexity},
i.e., the same procedures over a map-of-sets adjacency with per-query DFS,
obtained by disabling the closures with a switch; (b) plain Cypher
variable-length path queries expressing the same structural questions
declaratively; (c) Neo4j Graph Data Science traversal
procedures~\cite{neo4jgds} over a projected graph; and (d) round-tripping the
DAG to the external libraries dagitty~\cite{textor2016dagitty},
pgmpy~\cite{ankan2024pgmpy}, causal-learn~\cite{zheng2024causallearn}, and
DoWhy~\cite{sharma2020dowhy}, reporting export, build, and query time
separately since the export is the overhead an in-DBMS design removes. Not
every baseline can answer every operator: Table~\ref{tab:coverage} states
which comparison is made on which operator, and a system that cannot express
an operator is absent from the corresponding figure rather than charged an
arbitrary cost.

\begin{table}[t]\centering
\caption{Which baseline can answer which operator. Cypher and GDS express
reachability and path patterns only; d-separation and identification are
compared against the DFS baseline and the external libraries ($^\ast$dagitty
and pgmpy).}
\label{tab:coverage}
\footnotesize
\setlength{\tabcolsep}{3pt}
\begin{tabular}{@{}lccccc@{}}
\toprule
Operator & Index & DFS & Cypher & GDS & Ext. \\
\midrule
reachability, path test & yes & yes & yes & yes & yes \\
path roles              & yes & yes & yes & no  & yes \\
d-separation            & yes & yes & no  & no  & yes \\
identification          & yes & yes & no  & no  & yes$^\ast$ \\
\bottomrule
\end{tabular}
\end{table}

\noindent\textbf{Metrics.} For RQ1 we report median and tail ($p_{99}$)
latency and sustained throughput per operator; for RQ2, build time and
resident memory against $|V|$ and density for both the closure and the
layer-pruned branch; for RQ3, the wall-clock of an incremental batch against
a full recompile across batch sizes $\{1, 10, 100, 1000\}$ and edit types
(insertion, deletion, \textsf{active} flip, \textsf{latent} flip, variable
addition and removal), the distribution of cone sizes, and the query latency
observed by a stream interleaving one edit every $k \in \{1, 10, 100\}$
queries under three maintenance policies (external re-export,
invalidate-and-recompile, patch); for RQ4, latency against $|V|$ on a log-log
scale; for RQ5, exact agreement of d-separation verdicts, canonical
adjustment sets, and the family of minimal adjustment sets against dagitty
and pgmpy on every dataset.

\subsection{Query performance (RQ1)}

\TODO{Results. Expected structure: per operator, index vs. DFS baseline on
all datasets (speed-up as a function of $|V|$ for tier (a) operators; the
masked vs. unmasked d-separation ablation for tier (b); canonical decision
vs. the level-wise enumeration of the earlier prototype for tier (c), which
does not finish on \textsc{andes} and beyond); then index vs. Cypher and GDS
on the operators of Table~\ref{tab:coverage}; then index vs. the external
round-trip with the export cost separated. State the headline numbers used in
the abstract and the introduction here and nowhere else.}

\subsection{Index cost and crossover (RQ2)}

\TODO{Build time and memory vs. $|V|$ and density; both branches; the
$|V|$ at which the closures exceed the configured budget and the latency
paid by the layer-pruned branch beyond it. Table with memory and build time
per dataset.}

\subsection{Incremental maintenance (RQ3)}

\TODO{Patch vs. recompile by batch size and edit type (Fig.~\ref{fig:incr});
the cone-size distribution explaining the deletion cost; the threshold at
which the fallback is triggered; the mixed-workload latency under the three
policies. This is the experiment that quantifies R4.}

\subsection{Scalability (RQ4)}

\TODO{Latency vs. $|V|$ on the synthetic grid (Fig.~\ref{fig:latency}), one
line per system; the observed slope for tier (a) operators (expected flat
up to the word count), for masked d-separation (expected sublinear on sparse
DAGs, linear on dense ones), and for identification.}

\subsection{Correctness (RQ5)}

\TODO{Agreement counts: d-separation verdicts, canonical sets, minimal-set
families, per dataset, against dagitty and pgmpy. Report any disagreement
and its cause; the expected outcome is exact agreement since the criterion
implemented is the one dagitty implements.}

\begin{figure}[t]\centering
  \fbox{\parbox{0.9\linewidth}{\centering \TODO{Fig.: query latency vs.\
  $|V|$ (log-log), one line per system (index, DFS baseline, Cypher, GDS,
  external round-trip), one panel per operator tier.}}}
  \caption{Query latency against graph size (RQ4).}
  \label{fig:latency}
\end{figure}

\begin{figure}[t]\centering
  \fbox{\parbox{0.9\linewidth}{\centering \TODO{Fig.: incremental patch vs.\
  full recompile, by batch size and edit type; mark the fallback threshold;
  second panel: query latency in the mixed workload under the three
  policies.}}}
  \caption{Incremental maintenance against recompilation (RQ3).}
  \label{fig:incr}
\end{figure}
